% META: {"id": "1NSI-TC-EVAL-A", "chapitre": "1NSI-TYPES-CONSTRUITS", "type_objet": "evaluation", "capacites": ["1NSI-TYPES-CONSTRUITS-C1", "1NSI-TYPES-CONSTRUITS-C2", "1NSI-TYPES-CONSTRUITS-C3", "1NSI-TYPES-CONSTRUITS-C4", "1NSI-TYPES-CONSTRUITS-C5"], "mode_creation": "adaptation", "sources_inspiration": ["corpus_nsi/03_progressions/supports/premiere/P04/P04_evaluation_types_construits.md"], "duree_min": 55, "status": "needs_review"}
\section*{Évaluation A --- Types construits}
\textit{Durée : 55 minutes. Calculatrice interdite. Documents interdits.}

\baremeIndicatif{Ex. 1 : 4 pts (analyser) ; Ex. 2 : 6 pts (programmer) ;
Ex. 3 : 5 pts (analyser, programmer) ; Ex. 4 : 5 pts (programmer)}

\bigskip

\textbf{Exercice 1} --- \textit{Lecture de programme} \hfill (4 points)

On considère le programme suivant :
\begin{python}
joueurs = [
    {"pseudo": "Ace", "score": 150},
    {"pseudo": "Bob", "score": 200},
    {"pseudo": "Cat", "score": 180},
]
meilleur = joueurs[0]
for j in joueurs:
    if j["score"] > meilleur["score"]:
        meilleur = j
print(meilleur["pseudo"])
print(meilleur["score"])
\end{python}

% BEGIN-TRACE
% joueurs = [{"pseudo": "Ace", "score": 150}, {"pseudo": "Bob", "score": 200}, {"pseudo": "Cat", "score": 180}]
% meilleur = joueurs[0]
% for j in joueurs:
%     if j["score"] > meilleur["score"]:
%         meilleur = j
% print(meilleur["pseudo"])
% print(meilleur["score"])
% EXPECTED
% Bob
% 200
% END-TRACE

\begin{enumerate}
  \item Donner, sans exécuter le programme, les deux valeurs affichées.
  \item Expliquer en une phrase le rôle de ce programme.
  \item Si on ajoute \lstinline{meilleur["score"] = 0} après la boucle,
        quelle est la valeur de \lstinline{joueurs[1]["score"]} ? Justifier.
\end{enumerate}

\bigskip

\textbf{Exercice 2} --- \textit{Programmation avec des tableaux} \hfill (6 points)

Un tournoi de e-sport enregistre les scores de chaque manche dans un tableau :

\begin{python}
scores = [150, 200, 180, 210, 170]
\end{python}

\begin{enumerate}
  \item Écrire une fonction \lstinline{au_dessus(scores, seuil)} qui renvoie le tableau
        des scores strictement supérieurs au seuil.
  \item Écrire une fonction \lstinline{moyenne(scores)} qui renvoie la moyenne des scores.
        La fonction doit lever \lstinline{ValueError} si le tableau est vide.
  \item Écrire une fonction \lstinline{classement(scores)} qui renvoie le tableau des
        tuples \lstinline{(rang, score)} triés par score décroissant.
        Exemple : \lstinline{classement([150, 200, 180])} renvoie
        \lstinline{[(1, 200), (2, 180), (3, 150)]}.
\end{enumerate}

% BEGIN-VERIFY
% def au_dessus(scores, seuil):
%     return [s for s in scores if s > seuil]
% assert au_dessus([150, 200, 180], 170) == [200, 180]
% assert au_dessus([100], 200) == []
% def moyenne(scores):
%     if not scores: raise ValueError
%     return sum(scores) / len(scores)
% assert moyenne([150, 200, 180]) == 510 / 3
% END-VERIFY

\bigskip

\textbf{Exercice 3} --- \textit{Grille et pixels} \hfill (5 points)

On représente une image en niveaux de gris par une grille (tableau de tableaux) :
\begin{python}
image = [
    [100, 200, 100],
    [200, 50,  200],
    [100, 200, 100],
]
\end{python}

\begin{enumerate}
  \item Combien de lignes et de colonnes contient cette image ?
  \item Quelle est la valeur du pixel central ? Donner l'expression Python.
  \item Écrire une fonction \lstinline{nb_sombres(image, seuil)} qui compte le nombre
        de pixels dont la valeur est inférieure ou égale au seuil.
\end{enumerate}

% BEGIN-VERIFY
% def nb_sombres(image, seuil):
%     c = 0
%     for ligne in image:
%         for pixel in ligne:
%             if pixel <= seuil:
%                 c += 1
%     return c
% image = [[100, 200, 100], [200, 50, 200], [100, 200, 100]]
% assert nb_sombres(image, 100) == 5
% assert nb_sombres(image, 50) == 1
% assert nb_sombres(image, 0) == 0
% END-VERIFY

\bigskip

\textbf{Exercice 4} --- \textit{Dictionnaire et mutabilité} \hfill (5 points)

On gère un carnet de contacts :
\begin{python}
contacts = {"Ali": "55123456", "Sara": "55234567"}
\end{python}

\begin{enumerate}
  \item Écrire une fonction \lstinline{chercher(contacts, nom)} qui renvoie le numéro
        du contact ou \lstinline{"Inconnu"} si le nom n'est pas dans le carnet.
  \item Écrire une fonction \lstinline{ajouter(contacts, nom, numero)} qui ajoute un
        contact au carnet. Préciser : cette fonction modifie-t-elle le dictionnaire
        passé en argument ?
  \item On exécute :
\begin{python}
carnet1 = {"Ali": "55123456"}
carnet2 = carnet1
carnet2["Sara"] = "55234567"
\end{python}
        Quelle est la valeur de \lstinline{carnet1} après ces trois lignes ? Justifier.
\end{enumerate}

% BEGIN-VERIFY
% def chercher(contacts, nom):
%     return contacts.get(nom, "Inconnu")
% assert chercher({"Ali": "55123456"}, "Ali") == "55123456"
% assert chercher({"Ali": "55123456"}, "Bob") == "Inconnu"
% END-VERIFY
